\chapter{Testare și validare}\label{ch:testare}
\pagestyle{fancy}

Acest capitol descrie metodologia de testare aplicată sistemului de control
FPGA, acoperind atât testarea individuală a fiecărui modul cât și testarea
integrată a lanțului complet de semnal. Se prezintă rezultatele măsurătorilor,
verificarea funcționalității și analiza utilizării resurselor FPGA.

%==========================================================================
\section{Metodologia de testare}
\label{sec:metodologie}

Strategia de testare urmează o abordare de tip \textit{bottom-up}: fiecare
modul VHDL este verificat individual înainte de a fi integrat în sistemul
complet. Această metodologie permite izolarea defectelor la nivelul modulului
unde au origine, simplificând depanarea.

Testarea s-a desfășurat pe două niveluri:

\begin{enumerate}
	\item \textbf{Testare per modul:} Fiecare modul este instanțiat izolat
	      în modulul de nivel superior, cu semnale de intrare controlate
	      (hardcodate sau furnizate de un generator de stimuli) și ieșiri
	      observate pe LED-uri, osciloscp sau terminal serial.

	\item \textbf{Testare de integrare:} Modulele sunt conectate conform
	      arhitecturii finale, iar sistemul complet este exercitat prin
	      scenarii end-to-end: cadre UART transmise de la PC-ul gazdă,
	      verificarea unghiurilor pe osciloscop la ieșirile PWM și
	      monitorizarea indicatoarelor LED.
\end{enumerate}

Instrumentele utilizate:
\begin{itemize}
	\item Osciloscop digital pentru măsurarea semnalelor PWM și SPI
	\item Terminal serial (PuTTY / Tera Term) la 115\,200~baud pentru
	      comunicarea UART
	\item LED-urile plăcii Basys3 pentru monitorizarea stărilor interne
	\item Vivado Hardware Manager pentru programarea și reprogramarea
	      rapidă a bitstreamului
\end{itemize}

%==========================================================================
\section{Testarea modulului PWM}
\label{sec:test-pwm}

Modulul \texttt{pwm\_gen} a fost primul testat deoarece generarea corectă a
semnalelor PWM este condiția fundamentală pentru controlul servomotoarelor.

\subsection{Metodă}

S-a instanțiat o singură instanță a modulului \texttt{pwm\_gen} cu unghiul
hardcodat la diferite valori de test. Ieșirea PWM a fost conectată la un
pin Pmod și măsurată cu osciloscopul.

\subsection{Cazuri de test}

\begin{table}[H]
\centering
\caption{Rezultatele testării modulului PWM}
\label{tab:pwm-test}
\begin{tabular}{@{}cccc@{}}
\toprule
\textbf{Unghi intrare} & \textbf{Lățime puls așteptată} & \textbf{Lățime puls măsurată} & \textbf{Eroare} \\
\midrule
0° & 1,000\,ms & 1,000\,ms & $<$10\,ns \\
90° & 1,500\,ms & 1,500\,ms & $<$10\,ns \\
180° & 2,001\,ms & 2,001\,ms & $<$10\,ns \\
45° & 1,250\,ms & 1,249\,ms & $<$10\,$\mu$s \\
135° & 1,750\,ms & 1,751\,ms & $<$10\,$\mu$s \\
\bottomrule
\end{tabular}
\end{table}

\subsection{Analiza rezultatelor}

Verificarea la 90° este deosebit de relevantă deoarece corespunde poziției
mediane a servomotorului:

\begin{equation}
\text{threshold}_{90°} = 100\,000 + 90 \times 556 = 100\,000 + 50\,040 = 150\,040
\end{equation}

La 100\,MHz, 150\,040~de cicluri corespund unei lățimi de puls de
$150\,040 \times 10\,\text{ns} = 1\,500\,400\,\text{ns} \approx 1,500\,\text{ms}$,
ceea ce este corect.

Eroarea minimă observată (sub un ciclu de ceas de 10\,ns) confirmă precizia
generatorului PWM hardware. Comparativ, o implementare software pe
microcontroler ar prezenta jitter de ordinul microsecundelor din cauza
întreruperilor și a schimbărilor de context.

Perioada semnalului PWM a fost măsurată la exact 20,000\,ms ($\pm$10\,ns),
confirmând funcționarea corectă a numărătorului pe 21~de biți.

%==========================================================================
\section{Testarea comunicației UART}
\label{sec:test-uart}

\subsection{Testul de loopback}

Pentru verificarea isolată a modulelor \texttt{uart\_rx} și \texttt{uart\_tx},
s-a implementat o configurație de loopback: fiecare octet recepționat este
retransmis imediat prin \texttt{uart\_tx}. Conectarea se realizează prin
cablul USB integrat pe placa Basys3, care include un convertor
USB-la-UART.

\subsection{Procedura de test}

\begin{enumerate}
	\item Se deschide PuTTY (sau Tera Term) cu parametrii: portul COM
	      corespunzător plăcii Basys3, 115\,200~baud, 8N1 (8~biți de date,
	      fără paritate, 1~bit de stop).
	\item Se activează opțiunea \textit{local echo} în terminalul serial
	      pentru a vizualiza caracterele transmise alături de cele
	      recepționate.
	\item Se tastează caractere individuale și secvențe de text. Fiecare
	      caracter transmis trebuie să apară de două ori pe ecran: o dată
	      de la echo-ul local și o dată de la loopback-ul FPGA.
	\item Se transmit secvențe rapide de caractere pentru a verifica
	      funcționarea la debit maxim.
\end{enumerate}

\subsection{Rezultate}

Toate caracterele au fost recepționate și retransmise corect, inclusiv
la debit maxim continuu. Nu s-au observat pierderi de caractere sau
erori de recepție. Testul confirmă:

\begin{itemize}
	\item Funcționarea corectă a sincronizatorului cu dublu flip-flop
	      (nicio eroare de metastabilitate detectată)
	\item Eșantionarea la mijlocul bitului funcționează corespunzător
	\item Calculul parametrilor de baud (CLKS\_PER\_BIT = 867) este corect
	\item Transmisia și recepția LSB-first sunt implementate conform
	      standardului
\end{itemize}

%==========================================================================
\section{Testarea protocolului de cadre}
\label{sec:test-frame}

\subsection{Verificarea recepției cadrelor}

Testarea modulului \texttt{frame\_rx} s-a realizat prin transmiterea de
cadre de 10~octeți de la PC prin intermediul unui script Python care
construiește cadre cu unghiuri cunoscute și calculează checksum-ul XOR
corect.

\subsubsection{Cazul 1: Cadru valid}

S-a transmis un cadru cu toți servo-urile la 90° (0x5A):
\begin{center}
\texttt{FF 5A 5A 5A 5A 5A 5A 5A 5A [checksum]}
\end{center}

Checksum = $\texttt{FF} \oplus \texttt{5A} \oplus \texttt{5A} \oplus
\texttt{5A} \oplus \texttt{5A} \oplus \texttt{5A} \oplus \texttt{5A} \oplus
\texttt{5A} \oplus \texttt{5A} = \texttt{A5}$.

Cadrul complet: \texttt{FF 5A 5A 5A 5A 5A 5A 5A 5A A5}.

Rezultat observat:
\begin{itemize}
	\item LED 2 (frame\_valid) pulsează scurt la recepția cadrului
	\item LED 3 (checksum\_err) rămâne stins
	\item LED-urile 8--15 afișează valoarea 0x5A (90) în binar
	\item Servomotorul conectat se poziționează la 90°
\end{itemize}

\subsubsection{Cazul 2: Cadru cu checksum eronat}

S-a transmis același cadru, dar cu octetul de checksum modificat intenționat
(un bit inversat):

\begin{center}
\texttt{FF 5A 5A 5A 5A 5A 5A 5A 5A A4} (checksum corect ar fi A5)
\end{center}

Rezultat observat:
\begin{itemize}
	\item LED 2 (frame\_valid) rămâne stins
	\item LED 3 (checksum\_err) pulsează scurt
	\item Pozițiile servomotoarelor nu se modifică (mailbox-ul nu primește
	      date noi)
\end{itemize}

\subsubsection{Cazul 3: Cadru incomplet (test timeout)}

S-au transmis doar primii 5~octeți ai unui cadru, apoi s-a așteptat mai
mult de 1\,ms fără a continua. Rezultat: modulul revine automat la starea
S\_WAIT\_DIR, iar la transmiterea unui cadru complet ulterior, acesta este
recepționat corect. Mecanismul de timeout previne blocarea permanentă a
parser-ului.

\subsection{Verificarea transmisiei cadrelor de răspuns}

S-a transmis un cadru cu direcția 0xFE (cerere de citire). FPGA-ul a
răspuns cu un cadru de 10~octeți conținând 0xFE + cele 8~unghiuri filtrate
curente + checksum XOR. Cadrul recepționat pe PC a fost verificat:
checksum-ul a fost corect, iar unghiurile au corespuns cu valorile afișate
pe LED-uri.

%==========================================================================
\section{Testarea SPI și ADC}
\label{sec:test-spi-adc}

\subsection{Verificarea interfeței SPI}

Testarea modulului \texttt{spi\_master} s-a realizat prin monitorizarea
semnalelor SPI pe osciloscop, cu sondele conectate la pinii Pmod JC
(SCLK, MOSI, MISO, CS).

S-au verificat:
\begin{itemize}
	\item \textbf{Polaritatea SCLK:} SCLK este inactiv la nivel '1'
	      (CPOL=1), confirmat vizual pe osciloscop
	\item \textbf{Fazarea datelor:} MOSI se modifică pe frontul descrescător
	      al SCLK, iar MISO este eșantionat pe frontul crescător (CPHA=1)
	\item \textbf{Frecvența SCLK:} Măsurată la 5\,MHz ($\pm$0,1\,MHz),
	      conform calcului 100\,MHz / 20
	\item \textbf{Ordinea biților:} MSB first, verificat prin compararea
	      secvenței de biți pe osciloscop cu valoarea transmisă
\end{itemize}

\subsection{Verificarea comunicării cu AD7124-8}

Primul test de comunicare a constat în citirea registrului de identificare
(ID register, adresa 0x05) al AD7124-8. Valoarea citită a fost comparată
cu valoarea specificată în fișa de catalog (0x14 pentru AD7124-8). Citirea
corectă a ID-ului confirmă:

\begin{itemize}
	\item Funcționarea corectă a masterului SPI
	\item Integritatea izolării SPI prin MAX14850
	\item Conexiunea electrică corectă între FPGA și ADC
\end{itemize}

\subsection{Verificarea scanării canalelor}

După inițializarea completă a ADC-ului, s-a verificat că toate cele
8~canale sunt scanate ciclic prin monitorizarea biților 3:0 ai octetului
de status. Secvența observată (0, 1, 2, 3, 4, 5, 6, 7, 0, 1, \ldots)
confirmă configurarea corectă a tuturor registrelor de canal.

LED-ul 0 (adc\_cfg\_done) se aprinde la finalizarea secvenței de
inițializare, iar LED-ul 1 (meas\_valid) pulsează periodic la fiecare
ciclu complet de citire a tuturor celor 8~canale.

%==========================================================================
\section{Testarea filtrului Kalman}
\label{sec:test-kalman}

Filtrul Kalman a fost testat prin trei categorii de teste: convergență,
răspuns la treapta și rejecția zgomotului.

\subsection{Testul de convergență}

\textbf{Obiectiv:} Verificarea că filtrul converge de la starea inițială
(incertitudine mare, $P_\text{init} = 10,0$) la o estimare stabilă.

\textbf{Procedură:} S-a furnizat o măsurătoare constantă de 90° pe toate
cele 8~canale. S-a monitorizat evoluția unghiului filtrat pe LED-urile
8--15 (care afișează unghiul servo 0).

\textbf{Rezultat:} După aproximativ 20~de cicluri ($20 \times 20\,\text{ms}
= 400\,\text{ms}$), unghiul filtrat s-a stabilizat la 90°. În primele
cicluri, valoarea a crescut rapid de la 0° (valoarea inițială) către 90°,
cu o tranziție exponențială caracteristică filtrului Kalman. Comportamentul
este conform cu teoria: incertitudinea inițială mare ($P_{00} = 10,0$)
face ca câștigul Kalman $K_0$ să fie aproape 1,0 în primele cicluri,
permițând filtrului să ,,creadă'' rapid măsurătoarea. Pe măsură ce
$P_{00}$ scade, câștigul se stabilizează la o valoare intermediară.

\subsection{Testul de răspuns la treaptă}

\textbf{Obiectiv:} Verificarea comportamentului filtrului la o schimbare
bruscă a unghiului (de la 0° la 180°).

\textbf{Procedură:} După stabilizarea filtrului la 0°, s-a transmis
un cadru UART cu unghiul 180°. S-a monitorizat evoluția ieșirii filtrate.

\textbf{Rezultat:} Filtrul a atins 90\% din valoarea finală (162°) în
aproximativ 10~cicluri (200\,ms) și s-a stabilizat complet la 180° în
aproximativ 15~cicluri (300\,ms). Comportamentul este un compromis între
răspunsul rapid (determinat de $R$ mic, adică încredere în măsurătoare)
și stabilitatea (determinată de $Q$ mic, adică încredere în predicție).
Componentele de viteză ($x_1$) au demonstrat o creștere tranzitorie în
timpul tranziției, confirmând că mecanismul de învățare a vitezei
funcționează corect (ecuația~\ref{eq:update-x1}).

\subsection{Testul de rejecție a zgomotului}

\textbf{Obiectiv:} Compararea ieșirii filtrate cu măsurătoarea brută
ADC pentru a evalua capacitatea de filtrare.

\textbf{Procedură:} S-a poziționat manual un servomotor la un unghi
fix și s-au citit simultan valoarea brută ADC (prin cadre de răspuns
0xFE) și valoarea filtrată (de pe LED-uri și din unghiurile PWM).

\textbf{Rezultat:} Valoarea brută ADC a prezentat fluctuații de
$\pm$2--3° datorită zgomotului analogic și cuantizării. Ieșirea filtrată
a prezentat fluctuații de maximum $\pm$0,5°, demonstrând o reducere a
zgomotului cu factor de 4--6$\times$. Servomotorul controlat de ieșirea
filtrată a fost vizibil mai stabil decât un servomotor controlat direct
de valoarea brută.

%==========================================================================
\section{Testarea modului demo}
\label{sec:test-demo}

\subsection{Procedură}

\begin{enumerate}
	\item Se comută sw0 în poziția ON pe placa Basys3.
	\item Se verifică LED 5 (demo mode active) --- trebuie aprins.
	\item Se observă secvențele de mișcare ale servomotoarelor.
	\item Se comută sw0 înapoi în poziția OFF pentru a verifica
	      dezactivarea.
\end{enumerate}

\subsection{Rezultate}

Toate secvențele demonstrative au fost executate corect:

\begin{itemize}
	\item \textbf{Deschidere/închidere pumn:} Toate cele 4~degete
	      (pinky, ring, middle, index) se deschid și se închid
	      simultan, cu mișcarea durând 1~secundă per tranziție.
	\item \textbf{Fluturare:} Încheietura rotește alternativ stânga
	      și dreapta, revenind la centru, cu o durată de 0,5~secunde
	      per direcție.
	\item \textbf{Flexii individuale:} Fiecare deget se închide
	      individual în ordine (pinky, ring, middle, index), cu
	      celelalte rămânând deschise.
	\item \textbf{Thumbs up:} Degetele se închid, iar degetul mare
	      rămâne extins --- gest vizual corect.
	\item \textbf{Pointing:} Degetul arătător rămâne extins, restul
	      se închid.
	\item \textbf{Flexia cotului:} Servomotorul de cot flexează și
	      extinde.
\end{itemize}

La dezactivarea sw0, secvența se oprește imediat, iar la reactivare
pornește de la începutul ciclului (cadrul-cheie 0). Tranziția între
modul demo și modul UART se realizează fără discontinuități datorită
filtrării Kalman care netezește orice schimbare bruscă de unghi.

%==========================================================================
\section{Testare de integrare}
\label{sec:test-integration}

\subsection{Lanțul complet de semnal}

Testarea de integrare a verificat funcționarea lanțului complet:
\begin{center}
Cadru UART $\rightarrow$ frame\_rx $\rightarrow$ mailbox $\rightarrow$
source\_mux $\rightarrow$ kalman\_engine $\rightarrow$ pwm\_gen $\rightarrow$
ADUM1400 $\rightarrow$ Servomotor MG996R
\end{center}

\textbf{Procedură:} S-au transmis cadre UART cu unghiuri variabile de la
PC, generând mișcări ale servomotoarelor. S-a verificat vizual că mișcarea
mâinii robotice corespunde unghiurilor comandate.

\subsection{Măsurarea latenței end-to-end}

Latența totală de la recepția unui cadru UART până la modificarea efectivă
a semnalului PWM a fost estimată:

\begin{table}[H]
\centering
\caption{Contribuțiile la latența end-to-end}
\label{tab:latency}
\begin{tabular}{@{}lr@{}}
\toprule
\textbf{Etapă} & \textbf{Durată} \\
\midrule
Recepție UART (10 octeți $\times$ 86,8\,$\mu$s) & 868\,$\mu$s \\
Parsare cadru (instantanee, pipeline) & $<$1\,$\mu$s \\
Depunere mailbox (instantanee) & $<$1\,$\mu$s \\
Așteptare tick 20\,ms (caz mediu) & 10\,ms \\
Execuție filtru Kalman (8 canale) & $\approx$14,6\,$\mu$s \\
Propagare la ieșire PWM & $<$1\,$\mu$s \\
\midrule
\textbf{Total (caz mediu)} & $\approx$\textbf{10,9\,ms} \\
\textbf{Total (caz cel mai defavorabil)} & $\approx$\textbf{20,9\,ms} \\
\bottomrule
\end{tabular}
\end{table}

Latența dominantă este așteptarea tick-ului de 20\,ms în mailbox. Aceasta
este inerentă arhitecturii servo (ciclul de 50\,Hz) și nu poate fi redusă
fără a modifica perioadă de bază. Latența de procesare propriu-zisă a
FPGA-ului (parsare + Kalman + PWM) este sub 15\,$\mu$s, cu trei ordine de
mărime mai mică decât ciclul servo.

\subsection{Verificarea indicatoarelor LED}

S-au verificat toate combinațiile de stări LED în scenariile de funcționare:

\begin{table}[H]
\centering
\caption{Starea LED-urilor în diferite scenarii de funcționare}
\label{tab:led-states}
\begin{tabular}{@{}lcccccccc@{}}
\toprule
\textbf{Scenariu} & \textbf{L0} & \textbf{L1} & \textbf{L2} & \textbf{L3} & \textbf{L4} & \textbf{L5} & \textbf{L6} & \textbf{L7} \\
\midrule
După pornire & ON & blink & --- & --- & blink & OFF & --- & blink \\
Cadru valid & ON & blink & blink & --- & blink & OFF & --- & blink \\
Eroare checksum & ON & blink & --- & blink & blink & OFF & --- & blink \\
Mod demo & ON & blink & --- & --- & blink & ON & --- & blink \\
Răspuns 0xFE & ON & blink & blink & --- & blink & OFF & blink & blink \\
\bottomrule
\end{tabular}
\end{table}

\textit{Notă:} ,,blink'' indică pulsuri scurte periodice, ,,---'' indică
absența activității pe acel LED. LED-ul 0 (ADC config done) rămâne
permanent aprins după inițializare, confirmând configurarea reușită.

%==========================================================================
\section{Utilizarea resurselor FPGA}
\label{sec:resource-utilization}

Tabelul~\ref{tab:resources} prezintă utilizarea resurselor FPGA după
sinteză și implementare pentru circuitul integrat Xilinx Artix-7
XC7A35T (prezent pe placa Basys3).

\begin{table}[H]
\centering
\caption{Utilizarea resurselor FPGA (Artix-7 XC7A35T)}
\label{tab:resources}
\begin{tabular}{@{}lrrr@{}}
\toprule
\textbf{Resursă} & \textbf{Disponibil} & \textbf{Utilizat} & \textbf{Procent} \\
\midrule
LUT-uri (Look-Up Tables) & 20\,800 & $\sim$3\,000 & 14\% \\
Flip-flop-uri             & 41\,600 & $\sim$2\,500 & 6\% \\
Block RAM (BRAM36)        & 50      & 2            & 4\% \\
DSP48E1                   & 90      & 2--4         & 4\% \\
\bottomrule
\end{tabular}
\end{table}

\subsection{Analiza utilizării}

Gradul de utilizare global al FPGA-ului este redus (sub 15\% pe orice
categorie de resursă), ceea ce oferă mai multe avantaje:

\begin{enumerate}
	\item \textbf{Marjă pentru extinderi:} Resursele rămase permit
	      adăugarea de module suplimentare fără a necesita migrarea la
	      un FPGA mai mare. De exemplu, un filtru Kalman cu 3~stări
	      (adăugarea accelerației) ar necesita doar 2--4 DSP48E1
	      suplimentare.

	\item \textbf{Performanță temporală:} Gradul scăzut de utilizare
	      a LUT-urilor facilitează rutarea cu căi de propagare scurte,
	      garantând respectarea constrângerilor de timing la 100\,MHz
	      fără dificultăți.

	\item \textbf{Consum energetic redus:} Resursele neutilizate nu
	      contribuie la consumul dinamic, menținând disiparea termică
	      la un nivel scăzut.
\end{enumerate}

\subsection{Distribuția resurselor per modul}

Principalii consumatori de resurse sunt:

\begin{itemize}
	\item \textbf{kalman\_engine:} Domină utilizarea LUT-urilor datorită
	      aritmeticii Q16.16 (multiplicări, adunări pe 32~de biți) și a
	      memoriei de stare (40 cuvinte $\times$ 32~biți). Blocurile
	      DSP48E1 sunt utilizate pentru multiplicările implementate de
	      funcția \texttt{mul\_q16} și pentru operațiile de diviziune.

	\item \textbf{ad7124\_ctrl:} Consumă o parte semnificativă din
	      flip-flop-uri datorită numărului mare de registre de stare
	      (secvența de inițializare, numărătoare, registre de date).

	\item \textbf{pwm\_gen (8 instanțe):} Fiecare instanță necesită un
	      numărător pe 21~de biți și un comparator, totalizând
	      $8 \times 21 = 168$ flip-flop-uri și logica de comparație
	      aferentă.

	\item \textbf{demo\_rom:} ROM-ul cu 21~de cadre-cheie este sintetizat
	      în LUT-uri (nu în Block RAM) deoarece dimensiunea sa mică nu
	      justifică utilizarea unui bloc BRAM dedicat.
\end{itemize}

Cele 2~blocuri Block RAM sunt utilizate de sintetizator pentru
implementarea memoriei de stare a filtrului Kalman (\texttt{state\_ram},
40 cuvinte $\times$ 32~biți) și pentru registrele interne ale
controlerului ADC.
